\documentclass[11pt]{article}
\usepackage{amsmath,amssymb,amsthm}
\title{Clean-born H stable packet at first mixed entry: exact weights and unpaid coordinates}\author{}\date{}
\begin{document}\maketitle
\paragraph{The stored packet being tested.}
Start with the original $H$ source of MT1 and its actual source-linear coefficient
under a clean formation step. File the whole chosen dimension-six projection and lower/current/contact terms
as prescribed in MT37--40. Retain the covariant quadratic and cubic parts as
separate coordinates of the complete stable packet, together with its analytic
tail; their decomposition is not deletion of its nonlocal quadratic/cubic parts. Denote the resulting completed stable packet by $R_j^H$, with canonical
source degree $d=2$, fixed clean birth radius $\alpha$, and actual saved graded size
$A_j$. The filed pieces remain separate state coordinates. This operational
specification refers to the actual clean coefficient; it introduces no alternate
fine observable or new contact prescription.

Let $e$ be its first geometrically touched mixed entry. Use its entire current
representative $F_{j,e}$ of LM42, including every proper difference already born,
not $H$ reevaluated on the original microscopic field. A completed carrier meeting
the modified component is the nonprotected case; PE12's untouched-row identity
cannot be used. The manuscript does not supply a verified selected-history
range certificate for this case. Consequently the calculations below distinguish
the established LAST regular bound from its still unverified first mixed use.

\paragraph{The established last-regular envelope, with all three weights.}
On its actual admitted common regular readout, RG31 (pp.755--756) gives
\[
 |F_{j,e}|\le C A_j\rho^2\left[
 r^s q_H^2+r^\beta\frac{g_e}{\alpha}q_H^3
       +r^2\frac{g_e^2}{\alpha^2}q_H^4\right],\qquad
 r=L^{j-e},\quad\rho=g_e/g_j.                         \tag{1}
\]
The symbol $\beta$ in the exponent $r^\beta$ is the saved cubic locality exponent,
not the Wilson action coefficient. Original integrated root counts and engineering
factors are already included in (1); there is no extra microscopic volume to insert.
The original mesh and test/anchor conventions remain in $A_j$ and its coefficient
norm. In particular the sharper cubic factor for $d=2$ is retained.

MT22--25 (pp.686--687) defines the relevant canonical control as
\[
 q_H=1+\|z_H\|+\max_{b\in H}|B_b|.                    \tag{2}
\]
Here $z_H$ includes all external determining, collar and pinned data with their
actual first-order norm. Every such input is scaled in the canonical family.
The control of current and fractional rows comes from the admitted common
response, not from the size of link logarithms alone. An arbitrary held-fixed
physical background is not bounded canonical $z_H$.

\paragraph{What the ORIGINAL reserve can actually pay.}
In the original pre-gauge-removal operation $T_0$, split the coordinates in (2)
into $H_{\rm paid}$ and $H_{\rm unpaid}$. The former must be literal current stored
variables whose unspent energy is part of the S59/S86 ledger; a previously
integrated coordinate is never restored. Suppose their actual positive forms give
\[
 \max_{H_{\rm paid}}|B_b|\le c_H E_\omega^{1/2},\qquad
 T_0 e^{\theta E_\omega}\le B_\omega.                 \tag{3}
\]
The constant $c_H$ retains the actual coercivity/coordinate comparison and counts;
it is not declared uniform merely by calling the domain a component. Set
$Q=1+\|z_H\|+\max_{H_{\rm unpaid}}|B_b|$ and let $Q_*$ be a valid bound on its
actual retained support. For each fixed integer $m$,
\[
 (T_0 q_H^{2m})^{1/2}
 \le C_m B_\omega^{1/2}
 \left[Q_*^m+c_H^m\left(\sup_{u\ge0}u^m e^{-\theta u}\right)^{1/2}\right].
                                                               \tag{4}
\]
This follows from $q_H\le Q_*+c_H\sqrt{E_\omega}$, the elementary sum inequality,
and $T_01\le B_\omega$. It uses one reserve even at $m=4$; squaring the quartic
tail costs degree eight, not a second rare mass. A fixed spatial carrier power
can be paid only from its saved spare locality. Unpaid reference energy and fixed
outer data are not integrated by (3).

Where the literal carrier is the common regular readout in (1), Minkowski gives
\[
 B_\omega^{-1/2}\|F_{j,e}\|_{L^2(T_0)}
 \le C A_j\rho^2\left[r^s H_2+
 r^\beta\frac{g_e}{\alpha}H_3+
 r^2\frac{g_e^2}{\alpha^2}H_4\right],                 \tag{5}
\]
where $H_m=Q_*^m+C_{m,\theta}c_H^m$.
This is the explicit energy bookkeeping needed for the outer history weighted carrying calculation, rather than an
assertion that all old-coordinate energies are available. Its extension to the
actual touched mixed readout is the missing use of (1).

\paragraph{Quantified loss when the record does not pay an input.}
If $Q_*,c_H$ are uniformly bounded, (5) preserves all three old age/coupling
weights, with no absolute inverse-coupling loss. If instead an unpaid canonical
input is bounded only by the original cutoff $P_e$, (5) retains the factors
$P_e^2,(g_e/\alpha)P_e^3,(g_e/\alpha)^2P_e^4$.
When $g_eP_e\le c\alpha$, each is at most a fixed constant times $P_e^2$.
For $P_e=(\log g_e^{-2})^p$, this fallback therefore loses
$(\log g_e^{-2})^{2p}$ at fixed age, rather than a fourth power with its coupling
factors silently omitted. It is a loss in the available estimate, not a lower
bound for the actual coefficient.

On the full physical core a permissible unpaid reference is only bounded by
$Q_*\lesssim\alpha/g_e$. The three contributions in (5) then have the common
absolute prefactor
\[
 A_j\rho^2(\alpha/g_e)^2
       =A_j(\alpha/g_j)^2,                            \tag{6}
\]
with age powers $r^s,r^\beta,r^2$ respectively. This is precisely the $d=2$
physical-core cost, not a uniform canonical bound. At $e-j=1$, strict age cannot
absorb an absolute $g_j^{-2}$ or an unbounded cutoff logarithm. The separately
scaled physical norm in RG6/MT40 legitimately retains (6); it must not be used
as the canonical $O(A_j)$ incoming norm.

\paragraph{The domain of the target in the outer history weighted carrying calculation.}
Its exact raw covariance and complex comparison inequalities hold.
Its boxed unweighted insertion target is a scalar CANONICAL target (or a target
for an already appropriately graded carrier). Imposing that same unweighted
amplitude uniformly over the full fixed physical domain is not justified. The
physical/canonical intersection must retain its distinct amplitude and radius
factors. This does not promote the canonical mixed-entry bound to a
proved result.

\paragraph{Raw normalization and constants that cancel.}
With the actual interaction estimate and unspent reserve established on this same
record, the outer history weighted carrying calculation converts (5) into a raw attachment bound by multiplying its
right side by $C g_e B_\omega$. There is one $B_\omega$ and the original source-zero
body remains. Normalizing an individual rare history would add an unproved ratio.
A scalar depending only on the fixed outer variable cancels from covariance;
one may subtract it when bounding that covariance, but the original scalar row
remains filed. The envelope loss above is not evidence such a scalar survives.
Removing an unpaid-background loss in the genuinely mixed parts requires an
actual response or cancellation estimate for the STORED packet, not just this
constant cancellation.

\paragraph{Exact unresolved clause.}
The missing model assertion is that the actual pre-removal selected integral of
this $F_{j,e}$ obeys (5) with its intended canonical grade and uniformly paid
$Q_*,c_H$, or an equivalent direct averaged bound avoiding pointwise (1).
The two concrete checks within that assertion are: the current representative
on the original selected support is the same common admitted readout used in
RG31 (LM47 is not supplied by the old unmarked range); and its variable/exterior
inputs are actually owned by the original reserve or uniformly canonical.
LM50--52 preserves real saved values under retention but does not prove that
its enlarged halo has the original local reserve/analytic kernel. No new halo,
new gauge quotient, or re-used failed-test credit has been introduced here.
Original current, selector-boundary and simultaneous-wall versions remain open.

\paragraph{Coordinate estimate and loss with age.}
On a retained relative-tree record the coordinate bound is
\[
 |B_b^{(R)}|^2\le C m_b\sum_p n_{bp}
       [\beta_R s(V_p)+\beta_R s(W_p)],
\]
where $B^{(R)}=g_R^{-1}\log(VW^{-1})$, $W$ is the retained reference,
and $n_{bp}$ are literal plaquette-filling multiplicities. The common tree gauge,
unitary telescoping and Cauchy--Schwarz give the proof. The
reference action remains in $Q_*$; only unused original energies may enter (3).
Incidence polynomials can be paid by a fixed unused record-count tilt if the
whole filling box already belongs to the record, not by inventing a halo.

This provides an additional quantitative diagnostic of the required scale match.
For a genuinely retained record at scale $R$, a coordinate converted to coupling
$g_e$ has factor $g_R/g_e$. Its finite-difference stencil of assigned current length
$h$ has a generic norm bound with another factor $h^{-1}$. Thus the elementary
fallback for THIS logarithmic coordinate/stencil has
\[
 c_H\lesssim C_{\rm incidence}(g_R/g_e)h^{-1}.
\]
This is not yet an identification of the full MT22 potential/current input; any
additional per-length conversion required there must also be retained. If a
record at $R=j$ genuinely remains (never restore an integrated one) and its
assigned stencil is $h=r=L^{j-e}$, substitution of this fallback into the three
terms of (5) gives respectively
\[
 r^{s-2},\qquad (g_j/\alpha)r^{\beta-3},\qquad
 (g_j/\alpha)^2r^{-2},                                 \tag{7}
\]
up to the displayed incidence and reserve constants. The relative coupling
factors cancel as shown; the missing powers of age do not. This is a loss of the
available generic derivative estimate, not a lower bound for the stored packet.
Current-unit stencils do not have this particular loss. The actual admitted
old-to-current response must therefore either identify only such current inputs,
or improve the generic fine-stencil bound before it can yield MT36. Polynomial
moment payment alone does not repair (7).

\paragraph{Verification.}
The cubic and quartic powers in (1) agree with RG31 for $d=2$.
Substituting $Q_*=\alpha/g_e$ gives (6) in all three terms; the second moment
requires precisely the degree-eight reserve. Source MT22 explicitly includes the
external first-order data, and MT25a explicitly distinguishes physical from
canonical size. No estimate above
beyond \cite{W} and \cite{ND} is used.
\begin{thebibliography}{9}
\bibitem{W} Wilson-source estimates, volumes 1--10 (\texttt{lattice\_theory/source\_estimates} in this collection), cited by source page.
\bibitem{ND} \emph{Nonlinear dispersion and unequal spatial weights}, \texttt{lattice\_theory/block\_and\_transfer/conditional\_estimates/nonlinear\_dispersion\_and\_unequal\_spatial\_weights.tex} in this collection.
\end{thebibliography}
\end{document}
